(1) 取向下为正，Newton方程为 \(m\dot v=mg_*-\beta v^2\)。终端速度要求总力为零，故 \(v_t=\sqrt{mg_*/\beta}\)；该稳定平衡来自阻力强度随速度增加，前提是物体下落时速度非负。

(2) 设 \(\tau=v_t/g_*=\sqrt{m/(\beta g_*)}\)，分离变量 \(dv/(1-v^2/v_t^2)=g_*dt\)。初速零给 \(v(t)=v_t\tanh(t/\tau)\)，再积分得 \(y(t)=v_t\tau\ln\cosh(t/\tau)\)；速度随时间单调趋于终端值。

(3) 乘 \(v\) 得 \(d(mv^2/2)/dt=mg_*v-\beta v^3\)，故热耗散率为 \(\beta v^3\ge0\)。下落到任意时刻时，浮力修正的重力势能释放为 \(mg_*y\)；累计散热 \(Q=mg_*y-mv^2/2\ge0\)，没有把热耗散错记为有用动能。

(4) 消去时间并利用 \(\tanh^2=1-\operatorname{sech}^2\)，得 \(v(H)^2=v_t^2[1-\exp(-2\beta H/m)]\)。\(H\) 很小时 \(v^2\sim2g_*H\)，大高度时趋 \(v_t^2\)。此模型只在可用二次阻力近似的流动Reynolds数范围有效，极低Re需改用线性Stokes阻力。
